{{دیوان
| عنوان = سُلَیمیٰ
| ماخذ = ویکی ماخذ
| ماخذ_لنک = https://ur.wikisource.org/wiki/%D8%B3%D9%8F%D9%84%D9%8E%DB%8C%D9%85%D9%B0%DB%8C
| مقابلہ = اقبال سائبر لائبریری (اقبال اکادمی پاکستان)
| مقابلہ_لنک = https://www.iqbalcyberlibrary.net
}}
<poem>
جس کی نمود دیکھی چشمِ ستارہ بیں نے
خورشید میں، قمر میں، تاروں کی انجمن میں

صُوفی نے جس کو دل کے ظُلمت کدے میں پایا
شاعر نے جس کو دیکھا قُدرت کے بانکپن میں

جس کی چمک ہے پیدا، جس کی مہک ہویدا
شبنم کے موتیوں میں، پھُولوں کے پیرہن میں

صحرا کو ہے بسایا جس نے سکُوت بن کر
ہنگامہ جس کے دم سے کاشانۂ چمن میں

ہر شے میں ہے نمایاں یوں تو جمال اس کا
آنکھوں میں ہے سُلَیمیٰ! تیری کمال اس کا
</poem>
